\clearpage
\chapter*{Fe de erratas de la exposición n.º 212}
\addcontentsline{toc}{chapter}{Fe de erratas de la exposición n.º 212}
\setcounter{footnote}{2}
\src{301}

\noindent\textbf{5. TDTE III: Preesquemas cociente.}

\begin{center}
[Seminario Bourbaki, t. 13, 1960/61, n.º 212, 20 pp.]
\end{center}

\noindent\textbf{Página 15, afirmación (i).} --- El hecho de que $Y=X/R$ sea
cuasiproyectivo sobre $S$ parece, por el momento, haberse demostrado solo
cuando $R$ proviene de una relación de equivalencia.

\noindent\textbf{Páginas 18 y 19.} --- Como señalamos al final de la exposición
siguiente, la conjetura considerada es definitivamente falsa. El «hecho
positivo» señalado en la Observación 8.1 parece haber sido demostrado
simultáneamente por diversos autores (NAGATA, ROSENLICHT, GROTHENDIECK,
\ldots).

Señalemos que la aplicación hecha en TDTE V de la teoría del paso al cociente
desarrollada aquí a la construcción de esquemas de Picard puede sustituirse
también por un uso adecuado de los esquemas de Hilbert (cf. SHMT). Como se
indicó en el § 8, la laguna más importante de la teoría presentada
\src{302}
aquí es la falta de un criterio de existencia de cocientes por una relación de
equivalencia no propia, como las relaciones de equivalencia procedentes de
ciertas acciones del grupo proyectivo. Un teorema importante en esta dirección
ha sido obtenido por MUMFORD\footnote{MUMFORD (David). --- An elementary
theorem in geometric invariant theory, \emph{Bull. Amer. Math. Soc.}, t. 67,
1961, pp. 483--486.}. Para un refinamiento de su resultado y diversas
aplicaciones a la teoría, véase SHMT.
